\part{Empfehlungen und Abschluss}
\chapterimage{Bilder_und_Co/abschluss.jpg} % Chapter heading image

% -------------------------
% CHAPTER 13: Empfehlungen zur Anwendung
% -------------------------
\chapter{Allgemeine Empfehlungen}\label{chap:empfehlungen}

Die folgenden Empfehlungen verdichten Befunde aus EASA-Leitlinien, 
ISO/PAS 8800 und der beispielbezogenen Analyse zu einem pragmatischen 
Handlungsrahmen. Sie sind bewusst modell-, tool- und domänenneutral formuliert 
und gelten unabhängig vom herangezogenen Beispiel der Random-Forest-gestützten 
Wetterklassifikation im Onboard-Radar. Das Beispiel dient der Veranschaulichung. 
Die Empfehlungen richten sich auf vergleichbare Assistenzfunktionen und sind als 
übergreifende Vorgehensweisen zu verstehen, die im jeweiligen Projektkontext weiter auszuarbeiten sind.




\section*{Empfehlungskatalog}

\subsection*{Vorbereitung \& Grundlagen}
Zunächst stehen die Grundlagen im Fokus: Einsatzgrenzen klären und die Datenbasis sichern.

\medskip
\begin{itemize}
  \item \textbf{Erst ODD, dann Daten:} 
        Den Einsatzbereich (ODD) mit Umgebungsbedingungen, Betriebsgrenzen, Schnittstellen, Annahmen und Autoritätsgrad eindeutig definieren. 
        Daraus gezielt Datenerhebung (repräsentative Szenarien) und erforderliche Nachweisbereiche ableiten. 

  \item \textbf{Datenbasis früh sicherstellen:} \footnote{\ref{sec:lernparadigmen} \nameref{sec:lernparadigmen}}
        Ohne ausreichende, belastbare Daten ist kein Training möglich und ohne Training keine KI. 
        Daten frühzeitig erheben, aufbereiten und versionieren.

\end{itemize}

\subsection*{Entwicklung \& Integration}
Sind ODD und Datenlage gesetzt, richtet sich der Blick auf bestehende Ansätze, 
deren gezielte Adaption und eine pragmatische Mensch-Maschine-Interaktion.

\medskip
\begin{itemize}

  \item \textbf{Das Rad nicht neu erfinden:} \footnote{\ref{sec:sop-ueberblick} \nameref{sec:sop-ueberblick}}
        KI-Lösungen für Radarsysteme (auch Flugzeugradarsysteme (Militär)) existieren bereits und sind erprobt. 
        Bestehende Ansätze als Ausgangsbasis nutzen statt neu zu entwickeln.

  \item \textbf{Human-in-the-Loop pragmatisch halten:} \footnote{\ref{sec:concept_paper} \nameref{sec:concept_paper}}
        Menschzentrierung ist vorgegeben und prägt den Scope. Eine pragmatische Ausgestaltung reduziert 
        Komplexität und erleichtert Nachweise. Beispiel: Zu jedem KI-Hinweis wird eine Wahrscheinlichkeitsangabe in \% angezeigt.

\end{itemize}


\subsection*{Nachweis \& Zulassung}
Im Anschluss wird die Absicherung adressiert – schrittweise, nachvollziehbar und zielgerichtet.

\medskip
\begin{itemize}
      
  \item \textbf{Schrittweise Reife aufbauen:}  \footnote{\ref{sec:bedenken_zertifizierung} \nameref{sec:bedenken_zertifizierung}}
        Klar präzisierte Regeln reduzieren Risiken durch enge Grenzen. 
        So kann die Entwicklung früher beginnen und es entstehen Praxiserfahrungen sowie Evidenz, 
        was dazu beitragen kann, das in der Literatur beschriebene Henne-Ei-Problem zu entschärfen.

  \item \textbf{Übernahme mit Augenmaß:} \footnote{\ref{sec:abdeck_easa_iso} \nameref{sec:abdeck_easa_iso}}
        Maßnahmen und Evidenzen aus etablierten Automotive-Standards können als Vorlage herangezogen werden. 
        Dabei sind kontextspezifische Anpassungen für die zivile Luftfahrt erforderlich, eine 1:1-Übernahme ist nicht sinnvoll, 
        insbesondere bei SOTIF-nahen Inhalten.
 
  \item \textbf{Grundlagen priorisieren, Komplexität beherrschen:} \footnote{\ref{sec:bedenken_zertifizierung} \nameref{sec:bedenken_zertifizierung}}
        Einfache, gut verstandene Lösungen wählen und Reproduzierbarkeit sicherstellen, 
        um Obsoleszenzrisiken bei wechselnden Tools und Bibliotheken zu reduzieren. 
        Komplexität niedrig halten und Annahmen offenlegen, damit die Absicherung trotz 
        geringer Betriebserfahrung beherrschbar bleibt.

\end{itemize}

\paragraph{Abschluss und Überleitung:}
Der Empfehlungskatalog priorisiert Nachvollziehbarkeit vor Tempo und setzt auf in der Praxis tragfähige Schrittfolgen. 
Zuerst werden Einsatzgrenzen (ODD) und Datenlage stabilisiert, danach werden Evidenz und HMI konkretisiert. 
Abschließend wird die Autorität maßvoll erhöht.

Zwei erprobte Grundsätze aus der Data Science bleiben leitend: \footnote{\ref{sec:praxis} \nameref{sec:praxis}}
(I) 80\,\% der Arbeit sind Datenarbeit (Sammeln, Bereinigen, Labeln, Versionieren) – wer hier investiert, reduziert spätere Risiken; (II) start
simple, baseline first – einfache, überprüfbare Lösungen schaffen robuste Ausgangspunkte und verhindern Fehloptimierungen.
Mit diesen Prinzipien wird KI nicht „per Dekret“ sicher, aber systematisch absicherbar.









% -------------------------
% CHAPTER 18: Fazit und Ausblick
% -------------------------
\chapter{Fazit und Ausblick}\label{chap:fazit-ausblick}

\section{Fazit}

Die Arbeit ordnet die in der Luftfahrt geforderten Absicherungsansätze systematisch ein, 
vergleicht sie mit heute etablierten Praktiken und leitet daraus eine belastbare 
Entscheidungsgrundlage ab. Kernbeiträge sind die strukturierte Zuordnung von Forderungen und 
Abdeckungen, die Verdichtung typischer Lücken und Risiken sowie ein praxisnaher Empfehlungskatalog 
für eine schrittweise, evidenzbasierte Einführung von KI-Funktionen im Radar-Umfeld. Damit liegt ein 
konsistenter Rahmen vor, um KI-Anwendungen in luftfahrtnahen Radarsystemen gezielt zu bewerten und abzusichern.

\medskip
Die Coverage-Analyse zeigt: ISO/PAS~8800 adressiert wesentliche Bausteine wie ODD-Führung, Daten- und 
Label-Governance, Robustheit, Verifikation/Validierung und Monitoring, erreicht im luftfahrtspezifischen 
Sinn jedoch häufig nur Teilabdeckung, die projektspezifisch zu präzisieren ist. Gleichwohl ergibt sich unter 
realistischen Annahmen eine tragfähige Brücke: Die in der Luftfahrt geforderten Konzepte lassen sich überwiegend 
zuordnen, identifizierte Lücken sind explizit adressierbar.

\medskip
Für das Beispiel einer Random-Forest-gestützten Wetterklassifikation im Onboard-Radar bestätigt die Bewertung: 
Unter klar definierten Einsatzgrenzen (Assistenzfunktion Level~1, kein Online-Lernen) und mit den skizzierten 
Evidenzen erscheint eine Zulassungsfähigkeit erreichbar. Offene Punkte betreffen primär die projektspezifische 
Ausgestaltung (u.\,a.\ Betriebs-Monitoring, Randfallpolitik, gewählte Operating Points), nicht die grundsätzliche 
Machbarkeit.

\medskip
Die regulatorische Einordnung bleibt eindeutig: KI-spezifische Guidance ergänzt bestehende Software- und 
Systemnachweise, ersetzt sie aber nicht. Nachweise skalieren weiterhin mit Hazard/DAL, Modelle werden offline 
trainiert und zur Freigabe eingefroren; ein „Nachweisrabatt“ ist ausgeschlossen. In Summe stützt die Evidenz, 
dass KI-basierte Assistenzfunktionen in Radarsystemen heute unter wohldefinierten Bedingungen und mit konsequenter 
Nachweisführung verantwortbar eingeführt werden können, während höher autorisierte Anwendungen zusätzliche Reife 
und Präzisierung erfordern. 


\section{Risikoeinordnung des Beispiels}

Ausgehend vom betrachteten Beispiel — Random-Forest-gestützte Wetterklassifikation als 
Onboard-Assistenzfunktion im zivilen Flugzeug, welche vor Inbetriebnahme 
eingefroren ist — lassen sich die wesentlichen Risiken klar abgrenzen. 
Zu den Mindestvoraussetzungen zählen eine belastbare Datenbasis, qualifizierte oder abgesicherte 
Werkzeuge sowie eine jederzeitige Übersteuerbarkeit durch den Menschen.

Die zentralen fachlichen Risiken liegen in ODD-Verfehlungen und Szenarienlücken, in der 
Qualität und Repräsentativität von Daten und Labels sowie in der Wahl geeigneter Operating 
Points zwischen Fehlalarmen und Verpassern. Hinzu kommen Robustheitseinbußen bei möglichen 
Konzept- und Daten-Drifts im Betrieb und menschliche Faktoren durch 
missverständliche Hinweise oder Scheingenauigkeit. Prozessseitig sind Toolchain-Disziplin 
und Auditierbarkeit relevant. Zusätzlich gilt es auch Security zu überprüfen.

Unter den genannten Annahmen sind diese Risiken beherrschbar. Wirksam sind eine sauber 
gefasste ODD mit konservativen Verhalten außerhalb des Einsatzbereichs, eine dokumentierte 
Datenführung mit klaren Kennzeichnungsregeln, risikobasierte Wahl der Arbeitspunkte pro 
Wetterphänomen, Abdeckungsnachweise, ein knappes und übersteuerbares HMI sowie 
Drift-Indikatoren und geregelte Re-Trainingszyklen außerhalb des Betriebs. 
Damit verbleiben die Hauptrisiken in der Repräsentativität der Daten, der Wahl geeigneter 
Arbeitspunkte und der Langzeitrobustheit. Die grundsätzliche Machbarkeit der Assistenzfunktion 
unter den definierten Randbedingungen gilt somit als bestätigt.



\section{Ausblick}
Kurz- bis mittelfristig empfiehlt sich ein schrittweises Vorgehen: Mit enger Funktion, klarer ODD-Disziplin, 
expliziter Human-in-the-Loop-Rolle und messbaren KPIs starten; Evidenzketten stabilisieren, Betriebsdaten strukturiert zurückführen 
und die Autorität erst nach belastbaren Erfahrungen erweitern. Der vorgelegte Empfehlungskatalog bietet 
hierfür eine praxistaugliche Leitplanke.
\medskip

Forschung und Standardisierung sollten etablierte Best-Practice-Verfahren gezielt
nachweisfähig machen und zugleich verbleibende generische Lücken schließen. Dazu zählen
insbesondere (I) Hardware- und Tool-Zulassungen, (II) ein geregelter Umgang mit 
Re-Training bereits freigegebener oder im Einsatz befindlicher Modelle sowie (III) robuste
Betriebsindikatoren für OOD/Qualität. Langfristig sind darüber hinaus die
Grundlagen höherer Automatisierungsgrade (z.\,B.\ L3: Autoritätsgrenzen, Updates, Oversight)
methodisch und regulativ zu präzisieren. 

\paragraph{Pragmatische Leitlinie}
Zwei Grundsätze aus der Data Science bleiben leitend: (1) Datenarbeit dominiert den Aufwand— typisch entfallen 
\(\approx 80\,\%\) auf Beschaffung, Bereinigung, Labeling und Governance. (2) kleine, klar definierte Schritte 
sind zielführender als große Sprünge. Enge Scopes mit klaren Metriken liefern verlässlichere Evidenz als breit angelegte Funktionen. 
Beides spricht für frühe Data-Readiness, konservative Betriebsregeln (Abstain/Degrade) und konsequentes 
Lifecycle-Management.

\section{Schlussbild und Antwort auf die Forschungsfrage}
Unter realistischen Rahmenbedingungen und mit disziplinierter Evidenzführung ist KI im Radarumfeld als
Level-1-Assistenz hinreichend sicher machbar. 
Die Forschungsfrage ob KI sicher genug für ein Flugzeug ist, lässt sich wie folgt beantworten:

Ja, für klar eingegrenzte, assistierende Funktionen mit
stabilem ODD, kein Online-Lernen im Betrieb, konservativem Verhalten (Abstain/Degrade),
menschlicher Entscheidungsverantwortung und einer durchgängigen Evidenzkette. 

In diesem Rahmen sind die typischen Nachweisziele (DAL-proportional) realistisch erreichbar, etwa für die hier
betrachtete Random-Forest-gestützte Wetterphänomen-Assistenz.

Im Vordergrund steht dabei nicht „mehr Methode“, sondern die konsequente Operationalisierung:
ODD präzisieren und einhalten, Daten- und Tool-Governance führen, Kalibrierung/OOD überwachen und Änderungen
kontrolliert freigeben. Außerhalb dieser Leitplanken (höhere Autorität, offene ODDs, adaptives Online-Lernen)
ist die Aussage nicht pauschal übertragbar und bleibt Gegenstand weiterer Forschung und Regelwerksentwicklung.
